\documentclass[11pt]{article}
\usepackage[T1]{fontenc}
\usepackage{lmodern}
\usepackage[margin=1in]{geometry}
\usepackage{amsmath,amssymb,amsthm,mathtools}
\usepackage{aliascnt}
\usepackage{booktabs,array,tabularx}
\usepackage{enumitem}
\usepackage{microtype}
\usepackage{needspace}
\usepackage[numbers,sort&compress]{natbib}
\usepackage[hidelinks]{hyperref}
\usepackage[nameinlink,noabbrev]{cleveref}
\crefname{appendix}{appendix}{appendices}
\Crefname{appendix}{Appendix}{Appendices}
\setlength{\emergencystretch}{2em}
\clubpenalty=10000
\widowpenalty=10000
\displaywidowpenalty=10000
\newcolumntype{Y}{>{\raggedright\arraybackslash}X}
\allowdisplaybreaks[1]
\setlength{\textfloatsep}{16pt plus 2pt minus 2pt}
\setlength{\floatsep}{12pt plus 2pt minus 2pt}
\numberwithin{equation}{section}
\setlist[enumerate]{leftmargin=2.1em,itemsep=0.35em}

\newtheorem{theorem}{Theorem}[section]
\newaliascnt{lemma}{theorem}
\newtheorem{lemma}[lemma]{Lemma}
\aliascntresetthe{lemma}
\newaliascnt{proposition}{theorem}
\newtheorem{proposition}[proposition]{Proposition}
\aliascntresetthe{proposition}
\newaliascnt{corollary}{theorem}
\newtheorem{corollary}[corollary]{Corollary}
\aliascntresetthe{corollary}
\theoremstyle{definition}
\newaliascnt{definition}{theorem}
\newtheorem{definition}[definition]{Definition}
\aliascntresetthe{definition}
\theoremstyle{remark}
\newaliascnt{remark}{theorem}
\newtheorem{remark}[remark]{Remark}
\aliascntresetthe{remark}
\crefname{theorem}{theorem}{theorems}
\Crefname{theorem}{Theorem}{Theorems}
\crefname{lemma}{lemma}{lemmas}
\Crefname{lemma}{Lemma}{Lemmas}
\crefname{proposition}{proposition}{propositions}
\Crefname{proposition}{Proposition}{Propositions}
\crefname{corollary}{corollary}{corollaries}
\Crefname{corollary}{Corollary}{Corollaries}
\crefname{definition}{definition}{definitions}
\Crefname{definition}{Definition}{Definitions}
\crefname{remark}{remark}{remarks}
\Crefname{remark}{Remark}{Remarks}

\newcommand{\R}{\mathbb{R}}
\newcommand{\N}{\mathbb{N}}
\newcommand{\Nzero}{\mathbb{N}_0}
\newcommand{\norm}[1]{\left\lVert #1\right\rVert}
\newcommand{\ip}[2]{\left\langle #1,#2\right\rangle}
\newcommand{\Fclass}{\mathcal{F}_{=}^{\infty}}
\newcommand{\DetVal}{\mathsf{DetVal}}
\newcommand{\Nval}{N_{\epsilon,\mathrm{det}}^{\mathrm{val}}}
\newcommand{\Lip}{\operatorname{Lip}}
\newcommand{\Span}{\operatorname{span}}
\newcommand{\E}{\mathbb{E}}
\newcommand{\Prob}{\mathbb{P}}

\title{\textbf{Dimension-Dependent Lower Bounds for\\
Deterministic Zeroth-Order Stationarity}}
\author{Colombo Team\\[0.35em]
\small\href{mailto:colombo.2026@outlook.com}{\texttt{colombo.2026@outlook.com}}}
\date{}
\hypersetup{
  pdftitle={Dimension-Dependent Lower Bounds for Deterministic Zeroth-Order Stationarity},
  pdfauthor={Colombo Team},
  pdfsubject={Deterministic exact-value oracle complexity of smooth nonconvex stationarity},
  pdfkeywords={zeroth-order optimization, nonconvex stationarity, exact value oracle, deterministic lower bounds}
}

\begin{document}
\maketitle

\begin{abstract}
We study the number of exact function values required by a deterministic
algorithm to find an $\epsilon$-stationary point of a lower-bounded smooth
function on $\R^d$. The gradient is globally $L$-Lipschitz, the initial gap
is exactly $\Delta$, and queries and outputs are unrestricted. Writing
$a=L\Delta/\epsilon^2$, we prove that, for $d\geq8$ and $a>1/2$,
\[
 cd\min\!\left\{a,\frac{d}{\log(ed)}\right\}
 \leq N_{\epsilon,\mathrm{det}}^{\mathrm{val}}
 \leq(d+1)\lceil4a\rceil,
\]
where $c>0$ is universal. Thus coordinate finite differences have the
optimal query order $\Theta(da)$ when $1/2<a\leq d/\log(ed)$.
In every dimension, zero queries suffice exactly when $a\leq1/2$;
otherwise at least $d+1$ values are necessary when the initial value is
charged. The lower bound combines a smooth nonconvex chain with directions
selected after blocks of scalar queries. Flat gates preserve every earlier
reply exactly, radial compression permits arbitrarily distant queries,
and a final direction defeats an output that need not have been queried.
The completed witness is coercive, belongs to $C^\infty$, and has precisely
the prescribed initial gap. This gives a direct deterministic proof of the
multiplicative dimension--accuracy cost, with an explicit parameter range
and a sharp zero-query threshold. The argument does not determine the
complexity at higher accuracy in a fixed dimension.
\end{abstract}

\paragraph{Keywords.}
Zeroth-order optimization; nonconvex stationarity; exact value oracle;
deterministic query complexity; resisting oracle.

\medskip
\begin{samepage}
\noindent\textbf{AI Usage.}
Nearly the entire research pipeline for this paper was carried out by
Colombo, powered by GPT-5.6 Sol. 
\par
\end{samepage}

\section{Introduction}
\label{sec:intro}

How many scalar observations are needed to find a point at which a smooth
nonconvex function has a small gradient? With an exact gradient oracle,
gradient descent gives the familiar $O(L\Delta/\epsilon^2)$ bound, where
$L$ bounds the gradient Lipschitz constant and $\Delta$ bounds the initial
objective gap. With function values alone, coordinate finite differences
give a deterministic bound of order $dL\Delta/\epsilon^2$ in the
nontrivial accuracy range. The factor $d$ is easy to explain for this
algorithm: estimating one gradient requires a collection of scalar
observations. Proving that an arbitrary deterministic value-based algorithm
must incur the same factor is a different question.

An unrestricted value-based method need not estimate gradients, follow a
descent path, or return a previously queried point. It may probe anywhere
in $\R^d$ and use every digit of each exact real answer. Consequently,
neither charging $d$ values for a gradient nor multiplying a first-order
lower bound by $d$ establishes a value-oracle lower bound. A valid
construction must answer every adaptive query consistently with one fixed
smooth function and must keep the algorithm's actual output nonstationary.

We give such a construction for deterministic algorithms. The function
class consists of lower-bounded $C^\infty$ objectives with globally
$L$-Lipschitz gradients and a prescribed initial gap
$f(x_0)-\inf f=\Delta$. Only these quantitative bounds are imposed:
there is no bound on higher derivatives, on the query radius, or on the
location of the output. Internal exact-real computation is unrestricted.

\subsection{Results and their interpretation}

Set $a=L\Delta/\epsilon^2$, and let $N$ denote the deterministic minimax
number of exact scalar value queries. Our main result is
\begin{equation}
 cd\min\!\left\{a,\frac{d}{\log(ed)}\right\}
 \leq N\leq(d+1)\lceil4a\rceil,
 \qquad d\geq8,\quad a>\frac12,
 \label{eq:intro-lower}
\end{equation}
for a universal constant $c>0$. In particular,
\begin{equation}
 N=\Theta\!\left(\frac{dL\Delta}{\epsilon^2}\right)
 \quad\text{if}\quad
 d\geq8,\qquad \frac12<\frac{L\Delta}{\epsilon^2}
                   \leq\frac{d}{\log(ed)}.
 \label{eq:intro-matching}
\end{equation}
The finite-difference upper bound is a standard baseline. Our proof
establishes its optimal order against all deterministic adaptive policies
in the stated range, including policies with unbounded queries and
unqueried outputs.

The behavior near the trivial accuracy threshold is also explicit. For
every $d\geq1$, zero queries suffice if and only if $a\leq1/2$. Above
this threshold, at least $d+1$ queries are necessary under the convention
that $f(x_0)$ is charged. This last statement already holds for shifted
quadratics with known Hessian $LI_d$. It identifies the abrupt transition
from no information being needed to a number of scalar observations
proportional to the dimension.

For $a>d/\log(ed)$, the lower bound in \eqref{eq:intro-lower} remains
valid, but it no longer matches the upper bound. The minimum records the
range of the construction; it does not assert that the actual complexity
saturates. All statements remain valid when the exact-gap condition is
relaxed to $f(x_0)-\inf f\leq\Delta$.

These conclusions give a self-contained exact-gap refinement of rates
established in related formulations. The contribution is a complete proof
that treats normalization, adaptive transcript consistency, and terminal
outputs together, supplemented by the precise threshold and quadratic
baseline. \Cref{sec:related} explains the overlap with existing bounds.

\subsection{Proof mechanism}

The proof combines established analytic ingredients with a value-level
transcript construction. We use the smooth chain and radial compression
of Carmon et al.~\cite{CDHS2020}, together with the delayed-direction
approach used in exact-value convex lower bounds
\cite{WuEtAl2026Convex,WuEtAl2026SC} and in the nonconvex setting of
\citet{PTJ2026}. The chain is unchanged; the argument
below verifies the properties needed for the present stationarity model.

The chain's flat gates make its value independent of a suffix of
coordinates whenever the current coordinate is exactly zero and all
later coordinates are small. We answer a block of $\Theta(d)$ queries
before selecting its associated direction. That direction is then chosen
orthogonal to the entire block and nearly orthogonal to every earlier
query. The value identity ensures that these later choices leave all
previous replies unchanged. Radial compression puts every query into a
bounded geometric region, and a quadratic term prevents stationary escape
points outside that region.

After the transcript fixes the output, one final direction is selected
orthogonal to it. This preserves all query values while enforcing a
nonzero gradient at the output. The initial gap is invariant under
rotations of the frame, so its exact normalization is fixed before the
first reply. A chain of length $n$ costs order $dn$ scalar queries; the
gap permits $n$ proportional to $a$, while the projection estimate permits
$n$ proportional to $d/\log(ed)$. These two requirements give the
minimum in \eqref{eq:intro-lower}.

\paragraph{Organization.}
\Cref{sec:related} compares the theorem with existing bounds under their
respective oracle and success conventions. \Cref{sec:model} gives the
formal model and main theorem. \Cref{sec:elementary} proves the upper bound
and threshold, \Cref{sec:family} establishes the analytic properties of the
hard family, and \Cref{sec:transcript} constructs its exact transcript.
\Cref{sec:lower} completes the scaling and parameter argument.
\Cref{sec:scope} discusses the remaining scope.

\section{Related work and comparison}
\label{sec:related}

\paragraph{Directly comparable value-oracle bounds.}
Randomized gradient-free methods based on smoothing and random directions
are developed by Nesterov and Spokoiny~\cite{NS2017} and Ghadimi and
Lan~\cite{GL2013}. In the noiseless specialization of the latter, a
sufficiently small smoothing radius gives $O(da)$ scalar evaluations for
$\E\norm{\nabla f(X)}^2\leq\epsilon^2$. The deterministic upper bound
proved here has the same order and certifies the returned point on every
admissible objective.

Peng, Tang, and Jia~\cite[Theorem~E.8]{PTJ2026} establish the same
deterministic order $d\min\{a,d/\log(ed)\}$ for the smooth class with
gap at most $\Delta$, for all $d\geq1$ and $a\geq1$. Their result also
allows gradient projections charged by rank. Their
Corollary~G.10 gives $C^\infty$ witnesses in a comparable matched range.
Our theorem treats the exact-gap class and the full transition at
$a=1/2$, without claiming a better asymptotic rate or broader oracle.

A recent preprint of Xia and Yuan~\cite[Theorem~2.3]{XY2026} proves a
randomized exact-value lower bound under an additional global
function-Lipschitz bound $M$. In our notation, put
\[
 \zeta:=\min\!\left\{a,\frac{M^2}{\epsilon^2}\right\}.
\]
Their bound is $\Omega(d\zeta)$ when
$\zeta\geq c_0^{-1}$ and $d\geq C\zeta\log(2+\zeta)$, for universal
constants. Choosing $M^2\geq L\Delta$ gives $\Omega(da)$ in a comparable
high-dimensional range. This randomized result is also stronger in its
algorithm class than the bound proved here. \Cref{tab:comparison}
records the success conventions and parameter conditions.

\begin{table}[!t]
\centering
\small
\renewcommand{\arraystretch}{1.2}
\setlength{\tabcolsep}{4pt}
\caption{Scalar value-query bounds for smooth nonconvex stationarity.
Here $a=L\Delta/\epsilon^2>1/2$ and
$\zeta=\min\{a,M^2/\epsilon^2\}$. Queries and outputs are unrestricted.
An upper gap bound is sufficient for the established results and for our
upper bound; our lower-bound witnesses have gap exactly $\Delta$.
The constants $c_0,C$ in the randomized lower-bound row are universal.}
\label{tab:comparison}
\begin{tabularx}{\linewidth}{@{}>{\raggedright\arraybackslash}p{0.20\linewidth}>{\raggedright\arraybackslash}p{0.29\linewidth}Y@{}}
\toprule
Work & Value queries & Success criterion and conditions\\
\midrule
Coordinate finite\newline differences\newline
(\Cref{prop:upper})
& $(d+1)\lceil4a\rceil$
& Deterministic; $\norm{\nabla f(X)}\leq\epsilon$ on every instance;
all $d\geq1$ and $a>1/2$.\\
\addlinespace[5pt]
Ghadimi--Lan~\cite{GL2013}\newline(noiseless case)
& $O(da)$
& Randomized; $\E\norm{\nabla f(X)}^2\leq\epsilon^2$;
all dimensions, with a sufficiently small smoothing radius.\\
\addlinespace[5pt]
Peng--Tang--Jia\newline\cite[Theorem~E.8]{PTJ2026}
& $\displaystyle\frac{d}{500}\min\!\left\{a,\frac{d}{\log(ed)}\right\}$
\newline(lower bound)
& Deterministic; success on every instance; $d\geq1$, $a\geq1$;
globally $L$-smooth, gap at most $\Delta$. Also holds with
gradient projections charged by rank.\\
\addlinespace[5pt]
Xia--Yuan\newline\cite[Theorem~2.3]{XY2026}
& $\Omega(d\zeta)$
& Randomized; worst-case expected cost for success probability $2/3$;
$\Lip(f)\leq M$, $\zeta\geq c_0^{-1}$,
$d\geq C\zeta\log(2+\zeta)$.\\
\addlinespace[5pt]
This paper\newline(\Cref{thm:main})
& $\displaystyle\Omega\!\left(d\min\!\left\{a,\frac{d}{\log(ed)}\right\}\right)$
& Deterministic; worst-case cap for success on every instance;
$d\geq8$, $a>1/2$; exact-gap $C^\infty$ witnesses.\\
\bottomrule
\end{tabularx}
\end{table}

\paragraph{Deterministic derivative-free algorithms.}
The monograph of Conn, Scheinberg, and Vicente~\cite{CSV2009} and the
survey of Larson, Menickelly, and Wild~\cite{LMW2019} describe
finite-difference, interpolation, direct-search, and trust-region methods.
Convergence and evaluation-complexity results for these frameworks include
\citet{CSV2009TR,GJV2016,Vicente2013,CGT2012};
\citet{DRRZ2025} survey recent direct-search theory.
These methods manage the cost and geometry of local sampling. Our lower
bound applies to deterministic policies beyond any particular local-model
or step-selection framework. Randomized polling offers a different
mechanism, allowing useful descent directions with controlled probability
\cite{GRVZ2015}.

\paragraph{Stationarity lower bounds and exact transcripts.}
Vavasis~\cite{Vavasis1993} studies algorithm-independent black-box
difficulty for local minimization. Carmon et al.~\cite{CDHS2020}
establish the $\Omega(a)$ smooth-stationarity benchmark in sufficiently
high dimension for more informative derivative oracles. Their construction
provides the chain used here, but a derivative-oracle lower bound alone
does not yield the additional scalar-query factor $d$.
For exact-value convex minimization, Kerger~\cite{Kerger2026} proves a
near-quadratic dimension lower bound for Lipschitz objectives on a ball.
Wu et al.~\cite{WuEtAl2026Convex,WuEtAl2026SC} use exact prefix shielding
and batched delayed rotations for smooth convex and strongly convex
objective suboptimality with bounded query locations. The shared
transcript mechanism is relevant to our proof, while the objective class,
success criterion, and treatment of unbounded queries differ.
Noisy-oracle lower bounds such as \citet{Shamir2013} concern a different
information restriction: here every discrepancy in an exact scalar value
is observable, however small it is.

\paragraph{Stronger regularity and low dimension.}
Bounds under additional derivative assumptions should be compared
separately. Cubic regularization~\cite{NP2006}, accelerated first-order
methods~\cite{CDHS2017}, and the lower bounds in \citet{CDHS2021}
depend on quantitative higher-order smoothness. Doikov and
Grapiglia~\cite{DG2025} obtain $O(d^{3/2}\epsilon^{-3/2})$
derivative-free evaluations when the gap and Hessian-Lipschitz constant
are fixed. Membership in $C^\infty$ alone does not supply such a uniform
Hessian-Lipschitz bound.
Dimension also matters independently of regularity: Hinder~\cite{Hinder2018}
trades polynomial dimension costs for improved accuracy dependence using
derivative information and stronger assumptions. In one dimension,
Chewi, Bubeck, and Salim~\cite{CBS2023} obtain logarithmic deterministic
complexity for a combined value-and-gradient oracle. Neither result is a
value-only upper bound for the class considered here.

\section{Oracle model and main results}
\label{sec:model}

\subsection{The objective class}

Fix $d\in\N$, $L>0$, $\Delta\geq0$, $\epsilon>0$, and a specified
$x_0\in\R^d$. All norms and inner products are Euclidean, and $\log$ is the
natural logarithm. Set
\begin{equation}
    a:=\frac{L\Delta}{\epsilon^2}.
    \label{eq:a}
\end{equation}

\begin{definition}[Exact-gap smooth class]
The admissible class is
\begin{equation}
\begin{split}
\Fclass(d,L,\Delta;x_0):=\bigl\{f\in C^\infty(\R^d):\;&
    \inf_{x\in\R^d}f(x)>-\infty,\\
&\norm{\nabla f(x)-\nabla f(y)}\leq L\norm{x-y}
       \quad\text{for all }x,y\in\R^d,\\
& f(x_0)-\inf_{x\in\R^d}f(x)=\Delta\bigr\}.
\end{split}
\label{eq:class}
\end{equation}
\end{definition}

Here $L$ is a smoothness upper bound and need not be the smallest such bound. The absolute value $f(x_0)$ and the infimum are
not public; adding a constant does not change admissibility.
No convexity, compactness, minimizer-attainment, or higher-order smoothness
bound is assumed. The class is nonempty: a constant function suffices when
$\Delta=0$, and otherwise
$f(x)=\frac{L}{2}\norm{x-x_\star}^2$ is admissible whenever
$\norm{x_\star-x_0}=\sqrt{2\Delta/L}$.

\subsection{Policies and minimax complexity}

One charged oracle call at $x\in\R^d$ returns the exact real number $f(x)$.
The parameters $d,L,\Delta,\epsilon,x_0$ are public; function values,
including $f(x_0)$, are not supplied for free.

\begin{definition}[Deterministic exact-value policy]
For $T\in\Nzero$, a padded cap-$T$ policy consists of total maps
\[
    \pi_t:\R^{t-1}\longrightarrow\R^d\quad (1\leq t\leq T),
    \qquad
    \pi_{\mathrm{out}}:\R^T\longrightarrow\R^d.
\]
On an objective $f$, its queries, replies, and output satisfy
\begin{equation}
    x_t=\pi_t(y_1,\ldots,y_{t-1}),\qquad y_t=f(x_t),\qquad
    X_{\mathrm{out}}=\pi_{\mathrm{out}}(y_1,\ldots,y_T).
    \label{eq:policy}
\end{equation}
The class of such policies is denoted by $\DetVal_T$.
\end{definition}

The maps may depend on the public parameters. They need not be continuous,
measurable, computable, or restricted to finite precision. Queries may be
arbitrarily large or repeated, and the output need not have been queried.
An early-stopping method with worst-case cap $T$ is included by appending
ignored dummy queries after stopping. For $T=0$, the output is fixed by the
public parameters alone.

\begin{definition}[Minimax query complexity]
Define
\begin{equation}
\begin{split}
\Nval(d,L,\Delta;x_0):=
\inf\bigl\{T\in\Nzero:\;&\exists A\in\DetVal_T\ \forall f\in
\Fclass(d,L,\Delta;x_0),\\
&\norm{\nabla f(X_{\mathrm{out}}^A(f))}\leq\epsilon\bigr\},
\end{split}
\label{eq:complexity}
\end{equation}
with value $+\infty$ if the set is empty.
\end{definition}

Translating the objective, all queries, and the output shows that this
quantity does not depend on $x_0$. We therefore write
$\Nval(d,L,\Delta)$. A lower bound has the quantifier order
$\forall A\,\exists f$: the objective may depend on the entire deterministic
policy, but it must be one fixed function when that policy is run against it.

\subsection{Main theorem}

\begin{theorem}[Deterministic exact-value complexity]
\label{thm:main}
There is a universal constant $c>0$ with the following properties.
\begin{enumerate}[label=(\roman*)]
\item For every $d\geq1$,
\begin{equation}
    \Nval(d,L,\Delta)=0
    \quad\Longleftrightarrow\quad a\leq\frac12.
    \label{eq:threshold}
\end{equation}
\item If $a>1/2$, then in every dimension
\begin{equation}
    d+1\leq\Nval(d,L,\Delta)
    \leq(d+1)\lceil4a\rceil
    \leq8da+2d.
    \label{eq:basic-bounds}
\end{equation}
In particular, $\Nval(d,L,\Delta)=\Theta(d)$ for
$1/2<a\leq1$, with constants uniform over that interval and over all $d$.
\item If $d\geq8$ and $a>1/2$, then
\begin{equation}
    \Nval(d,L,\Delta)
    \geq c d\min\!\left\{a,\frac{d}{\log(ed)}\right\}.
    \label{eq:main-lower}
\end{equation}
Consequently, whenever $d\geq8$ and
$1/2<a\leq d/\log(ed)$,
\begin{equation}
    cda\leq\Nval(d,L,\Delta)\leq12da,
    \qquad
    \Nval(d,L,\Delta)=\Theta\!\left(\frac{dL\Delta}{\epsilon^2}\right).
    \label{eq:matching}
\end{equation}
\end{enumerate}
More precisely, every integer budget strictly below either asserted lower
bound, and every policy with that budget, admit a single coercive objective
in \eqref{eq:class} whose exact transcript leads to an output with gradient
norm strictly larger than $\epsilon$.
\end{theorem}

An explicit admissible value of $c$ is given in \eqref{eq:c}.
The dimension-dependent lower bound is proved directly for exact values;
it is not a reduction that charges $d$ value calls for a gradient call.
The minimum in \eqref{eq:main-lower} describes the bound proved here, not an
assertion that the actual complexity saturates for larger $a$.

\begin{remark}[An upper bound on the initial gap]
\label{rem:gap-upper}
Every bound in \cref{thm:main} also holds for the class obtained by replacing
$f(x_0)-\inf f=\Delta$ with $f(x_0)-\inf f\leq\Delta$.
The upper-bound proof uses only this inequality, and the lower-bound
witnesses satisfy the stronger equality. Thus this extension does not
require a monotonicity argument for the exact-gap classes.
\end{remark}


\begin{table}[!htbp]
\centering
\small
\renewcommand{\arraystretch}{1.2}
\caption{Parameter ranges in \Cref{thm:main}. The constants implicit in
$\Theta$ are uniform over each displayed range.}
\label{tab:regimes}
\begin{tabularx}{\linewidth}{@{}p{0.13\linewidth}p{0.28\linewidth}Y@{}}
\toprule
Dimension & Normalized gap $a$ & Deterministic exact-value complexity\\
\midrule
$d\geq1$ & $0\leq a\leq1/2$ & Exactly $0$.\\
$d\geq1$ & $1/2<a\leq1$ & $d+1\leq N\leq4(d+1)$; hence $\Theta(d)$.\\
$d\geq8$ & $1/2<a\leq d/\log(ed)$ & $cda\leq N\leq12da$; hence $\Theta(da)$.\\
$d\geq8$ & $a>d/\log(ed)$ & Lower bound $\max\{d+1,cd^2/\log(ed)\}$;
upper bound $(d+1)\lceil4a\rceil$.\\
\bottomrule
\end{tabularx}
\end{table}

\section{The upper bound and the zero-query threshold}
\label{sec:elementary}

\subsection{The smoothness inequality}

\begin{lemma}[Descent and the initial gap]
\label{lem:descent}
If $\nabla f$ is globally $L$-Lipschitz, then for all $x,y\in\R^d$,
\begin{equation}
    f(y)\leq f(x)+\ip{\nabla f(x)}{y-x}
                       +\frac{L}{2}\norm{y-x}^2.
    \label{eq:descent}
\end{equation}
If, in addition, $f$ is lower bounded, then
\begin{equation}
    \norm{\nabla f(x)}^2\leq2L\bigl(f(x)-\inf f\bigr).
    \label{eq:gap-gradient}
\end{equation}
\end{lemma}

\begin{proof}
The fundamental theorem of calculus and Lipschitz continuity give
\begin{align*}
 f(y)-f(x)-\ip{\nabla f(x)}{y-x}
 &=\int_0^1\ip{\nabla f(x+t(y-x))-\nabla f(x)}{y-x}\,dt\\
 &\leq\int_0^1 Lt\norm{y-x}^2\,dt
 =\frac{L}{2}\norm{y-x}^2.
\end{align*}
Taking $y=x-\nabla f(x)/L$ yields
$f(y)\leq f(x)-\norm{\nabla f(x)}^2/(2L)$.
Since $\inf f\leq f(y)$, this proves \eqref{eq:gap-gradient}.
\end{proof}

At the specified starting point,
\begin{equation}
    \norm{\nabla f(x_0)}^2\leq2L\Delta=2a\epsilon^2.
    \label{eq:start}
\end{equation}
Returning $x_0$ without a query therefore succeeds when $a\leq1/2$.

\subsection{A value-only algorithm}

For $h>0$, define the forward-difference vector
\begin{equation}
    [g_h(x)]_i:=\frac{f(x+he_i)-f(x)}{h},
    \qquad 1\leq i\leq d.
    \label{eq:fd}
\end{equation}
Computing it uses one base value and $d$ displaced values.

\begin{lemma}[Finite-difference error]
\label{lem:fd}
For every $x\in\R^d$,
\begin{equation}
    \norm{g_h(x)-\nabla f(x)}\leq\frac{Lh\sqrt d}{2}.
    \label{eq:fd-error}
\end{equation}
\end{lemma}

\begin{proof}
For each coordinate,
\[
 [g_h(x)]_i-\partial_i f(x)
 =\frac1h\int_0^h\bigl(\partial_i f(x+te_i)-\partial_i f(x)\bigr)\,dt.
\]
Its absolute value is at most
$h^{-1}\int_0^h Lt\,dt=Lh/2$. Squaring these $d$ inequalities and summing
proves the claim.
\end{proof}

Assume $a>1/2$ and fix
\begin{equation}
    h:=\frac{\epsilon}{4L\sqrt d},\qquad
    \delta:=\frac{\epsilon}{8},\qquad K:=\lceil4a\rceil.
    \label{eq:upper-parameters}
\end{equation}
Starting at $x_0$, at each stage $k=0,\ldots,K-1$ compute $g_k=g_h(x_k)$.
If $\norm{g_k}\leq7\epsilon/8$, return $x_k$ and stop. Otherwise set
\begin{equation}
    x_{k+1}:=x_k-g_k/L.
    \label{eq:update}
\end{equation}
For completeness, on an arbitrary scalar history that does not trigger
stopping in the first $K$ stages, define the output to be $x_0$.
That fallback will never be used on an admissible objective. After stopping,
ignored queries at $x_0$ may be appended to obtain a padded cap-$K(d+1)$ policy.

\begin{proposition}[Deterministic upper bound]
\label{prop:upper}
The preceding algorithm finds an $\epsilon$-stationary point using at most
$(d+1)\lceil4a\rceil$ values when $a>1/2$.
\end{proposition}

\begin{proof}
\Cref{lem:fd} gives $\norm{g_k-\nabla f(x_k)}\leq\delta$.
Thus, whenever the stopping test holds,
\[
    \norm{\nabla f(x_k)}\leq\norm{g_k}+\delta\leq\epsilon.
\]
It remains to show that the test must hold at one of the $K$ stages.
For a stage at which it fails, write $e_k=g_k-\nabla f(x_k)$.
By \eqref{eq:descent} and \eqref{eq:update},
\begin{align}
 f(x_{k+1})
 &\leq f(x_k)-\frac{1}{2L}\norm{g_k}^2
                    +\frac{1}{L}\ip{e_k}{g_k}\notag\\
 &\leq f(x_k)-\frac1L\left(\frac{\norm{g_k}^2}{2}
                                    -\delta\norm{g_k}\right).
 \label{eq:inexact}
\end{align}
The function $r\mapsto r^2/2-\delta r$ is strictly increasing for $r>\delta$.
Since $\norm{g_k}>7\epsilon/8$ and $\delta=\epsilon/8$,
\begin{equation}
    f(x_{k+1})<f(x_k)-\frac{35\epsilon^2}{128L}
                    <f(x_k)-\frac{\epsilon^2}{4L}.
    \label{eq:successful-descent}
\end{equation}
If all $K$ stopping tests failed, telescoping would give
\[
    f(x_K)<f(x_0)-\frac{K\epsilon^2}{4L}
           \leq f(x_0)-\Delta=\inf f,
\]
a contradiction. This argument does not require querying $f(x_K)$.
Each stage costs at most $d+1$ values, and there are at most $K$ stages.
Finally,
\[
    (d+1)\lceil4a\rceil\leq2d(4a+1)=8da+2d.
\]
\end{proof}

\subsection{A dimension lower bound from two quadratics}

\begin{proposition}[Dimension cost above the zero-query threshold]
\label{prop:quadratic}
If $a>1/2$, every deterministic policy that succeeds uniformly on
\eqref{eq:class} uses at least $d+1$ exact-value queries in the worst case.
This lower bound already holds for quadratics with Hessian $LI_d$.
\end{proposition}

\begin{proof}
Fix a budget $T\leq d$ and a policy $A\in\DetVal_T$.
Generate its transcript recursively with replies
\begin{equation}
 r_t:=\Delta+\frac L2\norm{x_t-x_0}^2.
 \label{eq:quadratic-replies}
\end{equation}
Let $p$ be the resulting output. If $T\geq1$, put $\bar x=x_1$ and choose a
unit vector $v$ orthogonal to
$\Span\{x_t-x_1:2\leq t\leq T\}$, whose dimension is at most
$T-1<d$. If $T=0$, put $\bar x=x_0$ and choose any unit vector $v$.
Set $r=\sqrt{2\Delta/L}$ and $\beta=Lr\ip{v}{\bar x-x_0}$, and define
\begin{equation}
 f_\pm(x):=\frac L2\norm{x-x_0\mp rv}^2\pm\beta.
 \label{eq:two-quadratics}
\end{equation}
Both functions are coercive and $C^\infty$, have Hessian $LI_d$, and satisfy
$f_\pm(x_0)-\inf f_\pm=\Delta$. For every queried point,
\[
 f_\pm(x_t)=r_t\mp Lr\ip{v}{x_t-\bar x}=r_t.
\]
Induction through the policy maps shows that both fixed objectives reproduce
the transcript and return the same point $p$. Their gradients obey
\[
 \frac{\norm{\nabla f_+(p)}^2+\norm{\nabla f_-(p)}^2}{2}
 =L^2\norm{p-x_0}^2+2L\Delta>\epsilon^2.
\]
Thus at least one objective has gradient norm strictly above $\epsilon$ at
the actual output, whether or not that point was queried.
\end{proof}

\begin{corollary}[Exact complexity for known-Hessian quadratics]
\label{cor:quadratic-exact}
On the subclass $f(x)=(L/2)\norm{x-c}^2+C$ with
$\norm{x_0-c}^2=2\Delta/L$ and unrestricted $C\in\R$, the deterministic
exact-value complexity is zero when $a\leq1/2$ and exactly $d+1$ when
$a>1/2$.
\end{corollary}

\begin{proof}
The zero-query guarantee follows from \eqref{eq:start}, and the lower bound
follows from \cref{prop:quadratic}. For the matching upper bound, fix any
$h>0$ and query $f(x_0)$ and $f(x_0+he_i)$ for $1\leq i\leq d$.
The vector with coordinates
\[
 g_i=\frac{f(x_0+he_i)-f(x_0)}h-\frac{Lh}{2}
\]
equals $\nabla f(x_0)=L(x_0-c)$ exactly. Return $x_0-g/L=c$.
\end{proof}

The additive shifts in \eqref{eq:two-quadratics} account for the convention
that the base value is charged: with $T$ queries, only $T-1$ independent
query differences need to be annihilated. Choosing between the two signs
handles an arbitrary output without imposing another orthogonality
constraint. In particular, \cref{prop:quadratic} proves the converse in
\eqref{eq:threshold}, including $d=1$, and gives
\begin{equation}
 d+1\leq\Nval(d,L,\Delta)\leq4(d+1)\leq8d
 \qquad\text{when }\frac12<a\leq1.
 \label{eq:transition}
\end{equation}

\section{A smooth hard family with an exact gradient certificate}
\label{sec:family}

The rest of the argument proves the dimension-dependent lower bound.
The construction must satisfy two competing requirements: it must hide
unselected directions exactly, yet keep a nonzero gradient at any output
missing the final direction. We first verify these analytic properties
for every orientation; the next section will choose an orientation without
changing any issued value.

\subsection{The smooth zero-chain}

We use the chain of \citet[Section~4]{CDHS2020}. Define
\begin{equation}
 \Psi(t):=
 \begin{cases}
 0,&t\leq\tfrac12,\\[1mm]
 \exp\!\left(1-\dfrac{1}{(2t-1)^2}\right),&t>\tfrac12,
 \end{cases}
 \qquad
 \Phi(t):=\sqrt e\int_{-\infty}^{t}e^{-s^2/2}\,ds.
 \label{eq:psi-phi}
\end{equation}
The zero extension of $\Psi$ is $C^\infty$: every right derivative at
$1/2$ is a polynomial in $(2t-1)^{-1}$ times
$\exp(-(2t-1)^{-2})$ and tends to zero there. Both functions and their first
two derivatives are bounded. For example, on $t>1/2$, writing $r=2t-1$,
\[
    \Psi'(t)=4r^{-3}\Psi(t),\qquad
    \Psi''(t)=(16r^{-6}-24r^{-4})\Psi(t),
\]
which are bounded near $r=0$ and tend to zero as $r\to\infty$.
For $\Phi$, the derivative formulas are
$\Phi'(t)=\sqrt e\,e^{-t^2/2}$ and
$\Phi''(t)=-t\sqrt e\,e^{-t^2/2}$.
Set
\begin{equation}
    P_j:=\norm{\Psi^{(j)}}_\infty,\qquad
    Q_j:=\norm{\Phi^{(j)}}_\infty,\qquad j=0,1,2.
    \label{eq:PQ}
\end{equation}
All these constants are finite. We will use
\begin{equation}
 \Psi,\Psi',\Phi,\Phi'\geq0,\qquad
 \Psi(1)=1,\quad \Psi(t)\geq1\ (t\geq1),\qquad
 \Phi'(t)>1\ (|t|<1).
 \label{eq:properties}
\end{equation}

For $n\geq1$, define $h_n:\R^n\to\R$ by
\begin{equation}
    h_n(z):=-\Psi(1)\Phi(z_1)
      +\sum_{i=2}^{n}\bigl[
       \Psi(-z_{i-1})\Phi(-z_i)-\Psi(z_{i-1})\Phi(z_i)\bigr].
    \label{eq:chain}
\end{equation}
The convention $z_0=1$ will be used only in derivative formulas.

\begin{lemma}[Chain estimates and exact value shielding]
\label{lem:chain}
Define the universal constants
\begin{equation}
\begin{split}
    C_h&:=4P_0Q_0,\\
    G&:=\max\{1,2P_0Q_1+2P_1Q_0\},\\
    H&:=2P_0Q_2+2P_2Q_0+4P_1Q_1.
\end{split}
\label{eq:chain-constants}
\end{equation}
For every $n\geq1$, the following statements hold.
\begin{enumerate}[label=(\roman*)]
\item The range diameter of $h_n$ is at most $C_hn$.
\item For every $z\in\R^n$,
\begin{equation}
    \norm{\nabla h_n(z)}\leq G\sqrt n,
    \qquad \norm{\nabla^2h_n(z)}_{\mathrm{op}}\leq H.
    \label{eq:chain-bounds}
\end{equation}
\item If $|z_i|<1$ for some $i\leq n$, then there is $j\leq i$ such that
$|z_j|<1$ and $\partial_jh_n(z)<-1$.
\item If $z_j=0$ and $|z_k|<1/2$ for all $k>j$, then
\begin{equation}
    h_n(z)=h_n(z_1,\ldots,z_{j-1},0,\ldots,0).
    \label{eq:shield}
\end{equation}
\end{enumerate}
\end{lemma}

\begin{proof}
Each of the $n-1$ links in the sum in \eqref{eq:chain} has absolute value
at most $2P_0Q_0$, and the first term has absolute value at most $P_0Q_0$.
Thus $|h_n(z)|\leq(2n-1)P_0Q_0\leq2nP_0Q_0$. The range diameter is at most
$4nP_0Q_0=C_hn$.

A coordinate occurs only in the link ending at that coordinate and in the
next link. Differentiating these two links gives, respectively, an absolute
bound $2P_0Q_1$ and $2P_1Q_0$. The first term and the last coordinate satisfy
the same bounds after the absent terms are omitted. Hence
$|\partial_jh_n(z)|\leq G$ for every $j$, proving the gradient bound.

The Hessian is tridiagonal. The absolute value of a diagonal entry is at
most $2P_0Q_2+2P_2Q_0$, and each adjacent off-diagonal entry has absolute
value at most $2P_1Q_1$. Every absolute row sum is therefore at most $H$.
For a real symmetric matrix the Euclidean operator norm is bounded by its
maximum absolute row sum, so the Hessian bound follows.

For part (iii), choose the smallest $j\leq i$ with $|z_j|<1$.
Then $|z_{j-1}|\geq1$, with $z_0=1$ at the first coordinate.
Differentiating \eqref{eq:chain} gives
\begin{align}
 \partial_jh_n(z)
 ={}&-\Psi(-z_{j-1})\Phi'(-z_j)-\Psi(z_{j-1})\Phi'(z_j)\notag\\
    &-\Psi'(-z_j)\Phi(-z_{j+1})-\Psi'(z_j)\Phi(z_{j+1}),
 \label{eq:chain-derivative}
\end{align}
where the second line is omitted if $j=n$. For $j=1$, the additional first
term involving $\Psi(-z_0)=\Psi(-1)$ is zero, so the formula remains valid.
Every term in \eqref{eq:chain-derivative} is nonpositive. At least one of
$\Psi(-z_{j-1})$ and $\Psi(z_{j-1})$ is at least one, and both
$\Phi'(-z_j)$ and $\Phi'(z_j)$ exceed one. Thus one of the first two terms
is strictly less than $-1$.

For part (iv), the link indexed by $j+1$ vanishes because
$\Psi(z_j)=\Psi(-z_j)=0$. Every later link, indexed by $i\geq j+2$, vanishes
because $|z_{i-1}|<1/2$. The first $j$ terms are unchanged when all
coordinates from $j$ onward are replaced by zero, since $z_j$ already is
zero. This proves \eqref{eq:shield}.
\end{proof}

The equality in \eqref{eq:shield} is exact, not an error estimate.
Small nonzero historical projections will only be used inside this
identity; no continuity of the algorithm's response to its observations is
required.

\subsection{Radial compression}

Fix the universal constants and radius
\begin{equation}
    \eta:=\frac1{100},\qquad
    D:=\frac{4(G+1)}{\eta},\qquad R:=D\sqrt n,
    \label{eq:radial-constants}
\end{equation}
and set
\begin{equation}
    \rho(x):=\frac{x}{\sqrt{1+\norm{x}^2/R^2}},
    \qquad x\in\R^d.
    \label{eq:rho}
\end{equation}
As in \citet[Section~5.2]{CDHS2020}, radial compression is combined with a
quadratic term. For an orthonormal frame
$U=(u_1,\ldots,u_n)\in\R^{d\times n}$, define
\begin{equation}
    F_U(x):=h_n(U^\top\rho(x))+\eta\norm{x}^2.
    \label{eq:FU}
\end{equation}
We give the derivative and gradient estimates for the constants above.

\begin{lemma}[Derivatives of the radial map]
\label{lem:rho}
The map $\rho$ is $C^\infty$, has $\norm{\rho(x)}<R$ for all $x$, and satisfies
\begin{equation}
    \norm{D\rho(x)}_{\mathrm{op}}\leq1,
    \qquad
    \norm{D^2\rho(x)[v,w]}\leq\frac6R
    \label{eq:rho-bounds}
\end{equation}
for all unit vectors $v,w$.
\end{lemma}

\begin{proof}
Let $\sigma=(1+\norm{x}^2/R^2)^{-1/2}$ and $y=\rho(x)=\sigma x$.
The denominator is everywhere positive, so smoothness is immediate.
Also $\norm{y}<R$, and differentiation gives
\begin{equation}
    D\rho(x)=\sigma I-\frac{\sigma^3}{R^2}xx^\top
             =\sigma\left(I-\frac{yy^\top}{R^2}\right).
    \label{eq:Drho}
\end{equation}
The matrix in parentheses has eigenvalues in $(0,1]$, which proves the
first bound. A second differentiation yields
\begin{align}
 D^2\rho(x)[v,w]
 ={}&-\frac{\sigma^3}{R^2}
       \bigl(\ip{x}{w}v+\ip{x}{v}w+\ip{v}{w}x\bigr)\notag\\
    &+\frac{3\sigma^5}{R^4}\ip{x}{v}\ip{x}{w}x.
 \label{eq:D2rho}
\end{align}
For $r=\norm{x}$ and $u=r/R$,
\[
    \frac{\sigma^3r}{R^2}
     =\frac1R\frac{u}{(1+u^2)^{3/2}}\leq\frac1R,
    \qquad
    \frac{3\sigma^5r^3}{R^4}
     =\frac3R\frac{u^3}{(1+u^2)^{5/2}}\leq\frac3R.
\]
The three terms in the first line of \eqref{eq:D2rho} contribute at most
$3/R$ in total, and the second line contributes at most $3/R$.
\end{proof}

\subsection{Smoothness and the exact gap}

The oracle replies must use their final scale from the beginning. The next
lemma shows that this is possible: changing the frame preserves the initial
gap, while one smoothness bound controls every orientation.

\begin{lemma}[A uniform hard family]
\label{lem:family}
For every orthonormal $d\times n$ frame $U$, the function $F_U$ is $C^\infty$
and coercive. Define
\begin{equation}
    \Lambda:=H+\frac{6G}{D}+2\eta.
    \label{eq:Lambda}
\end{equation}
Then $\Lip(\nabla F_U)\leq\Lambda$.
Moreover, the gap
\begin{equation}
    \delta_{n,d}:=F_U(0)-\inf_{x\in\R^d}F_U(x)
    \label{eq:gap-nd}
\end{equation}
is independent of the orientation of $U$ and satisfies
\begin{equation}
    \frac{e}{2\Lambda}\leq\delta_{n,d}\leq C_hn.
    \label{eq:gap-bounds}
\end{equation}
\end{lemma}

\begin{proof}
For unit vectors $v,w$, the second derivative of the chain term in
\eqref{eq:FU} is
\begin{align*}
 D^2[h_n(U^\top\rho(\cdot))](x)[v,w]
 ={}&\ip{\nabla^2h_n(U^\top\rho(x))U^\top D\rho(x)v}
              {U^\top D\rho(x)w}\\
 &+\ip{U\nabla h_n(U^\top\rho(x))}{D^2\rho(x)[v,w]}.
\end{align*}
Since $U$ is an isometry on $\R^n$, \cref{lem:chain,lem:rho} bound the
absolute value by $H+6G\sqrt n/R=H+6G/D$.
Adding the quadratic Hessian $2\eta I$ gives
$\norm{\nabla^2F_U(x)}_{\mathrm{op}}\leq\Lambda$.
Integrating the Hessian on a segment proves the claimed global Lipschitz
bound for the gradient.

The chain is bounded, whereas $\eta\norm{x}^2\to\infty$ as
$\norm{x}\to\infty$. Thus $F_U$ is coercive, lower bounded, and attains its
minimum. Since the quadratic is nonnegative,
\[
    F_U(0)-\inf F_U\leq h_n(0)-\inf h_n\leq C_hn.
\]
At the origin, \eqref{eq:chain-derivative} gives
$\nabla h_n(0)=-\sqrt e\,e_1$. As $D\rho(0)=I$,
\begin{equation}
    \nabla F_U(0)=-\sqrt e\,u_1.
    \label{eq:gradient-origin}
\end{equation}
Applying \cref{lem:descent} with smoothness constant $\Lambda$ gives
$\delta_{n,d}\geq e/(2\Lambda)$.

Finally, if $V$ is another orthonormal $d\times n$ frame, there exists an
orthogonal $d\times d$ matrix $Q$ with $V=QU$. The radial identity
$\rho(Qx)=Q\rho(x)$ implies
\[
    F_V(Qx)=h_n(V^\top\rho(Qx))+\eta\norm{Qx}^2=F_U(x).
\]
The map $x\mapsto Qx$ is a bijection and fixes the origin, so both the
infimum and the value at the origin are the same for the two frames.
This proves orientation independence.
\end{proof}

\subsection{A certificate at an arbitrary output}

Compression controls distant queries but can attenuate the chain gradient.
The following result separates the region where that gradient survives
from the region where the quadratic term dominates. This gives a global
certificate at the eventual output.

\begin{lemma}[Global final-direction certificate]
\label{lem:certificate}
Let $\gamma:=3/4$. For any orthonormal frame $U$ and any $x\in\R^d$,
\begin{equation}
    \ip{u_n}{\rho(x)}=0
    \quad\Longrightarrow\quad
    \norm{\nabla F_U(x)}>\gamma.
    \label{eq:certificate}
\end{equation}
\end{lemma}

\begin{proof}
Set $y=\rho(x)$, $z=U^\top y$, $w=U\nabla h_n(z)$, and
$\sigma=(1+\norm{x}^2/R^2)^{-1/2}$. The hypothesis is $z_n=0$.
By \cref{lem:chain}(iii), there exists $j\leq n$ such that
\begin{equation}
    |z_j|<1,\qquad
    \alpha_j:=\ip{u_j}{w}=\partial_jh_n(z)<-1.
    \label{eq:bad-coordinate}
\end{equation}
Since $D\rho(x)$ is symmetric, the chain rule and \eqref{eq:Drho} give
\begin{equation}
    \ip{u_j}{\nabla F_U(x)}
    =\sigma\left(\alpha_j-
           \frac{\ip{u_j}{y}\ip{y}{w}}{R^2}\right)
       +2\eta\ip{u_j}{x}.
    \label{eq:coordinate-FU}
\end{equation}

First suppose $\norm{x}\leq R/2$. Then
$\sigma\geq2/\sqrt5>4/5$ and $1/\sigma<5/4$.
Using $|z_j|<1$, $\norm{y}<R$, and $\norm{w}\leq G\sqrt n$,
\[
    \left|\frac{\ip{u_j}{y}\ip{y}{w}}{R^2}\right|
       <\frac{G\sqrt n}{R}=\frac GD<\frac1{400}.
\]
Also $x=y/\sigma$ implies $|\ip{u_j}{x}|<1/\sigma$.
Substitution into \eqref{eq:coordinate-FU} therefore gives
\begin{align*}
    \ip{u_j}{\nabla F_U(x)}
    &<-\sigma(1-G/D)+2\eta/\sigma\\
    &<-\frac45\left(1-\frac1{400}\right)+\frac{5\eta}{2}
      =-\frac{773}{1000}<-\frac34.
\end{align*}
Hence the gradient norm exceeds $3/4$.

If instead $\norm{x}\geq R/2$, the quadratic dominates globally:
\begin{align*}
    \norm{\nabla F_U(x)}
    &\geq2\eta\norm{x}-\norm{D\rho(x)^\top w}\\
    &\geq\eta R-G\sqrt n
      =(\eta D-G)\sqrt n
      =(3G+4)\sqrt n>\frac34.
\end{align*}
The two cases cover every point of $\R^d$.
\end{proof}

\section{Delayed directions and exact transcript completion}
\label{sec:transcript}

The analytic family is now fixed, but its frame $U$ is still free. We use
that freedom to complete a single objective consistent with an arbitrary
deterministic policy. Exact orthogonality hides a new direction from its
current block; the cap estimate below keeps its effects on older blocks
inside the chain's flat gates. The estimate is applied only after those
points have been fixed.

\subsection{Avoiding finitely many spherical caps}

\begin{lemma}[Simultaneous small projections]
\label{lem:caps}
Let $W\subseteq\R^d$ be a subspace of dimension $k\geq2$, and let
$z_1,\ldots,z_M\in\R^d$ satisfy $\norm{z_i}\leq R$, where
$M\in\Nzero$. For $0<\tau<R$, if
\begin{equation}
    2M\exp\!\left(-\frac{(k-1)\tau^2}{2R^2}\right)<1,
    \label{eq:cap-condition}
\end{equation}
then there is a unit vector $u\in W$ such that
\begin{equation}
    |\ip{u}{z_i}|<\tau\qquad (1\leq i\leq M).
    \label{eq:cap-conclusion}
\end{equation}
\end{lemma}

\begin{proof}
For $M=0$, any unit vector in $W$ suffices. Otherwise let $u$ be uniform
on the unit sphere of $W$. We first prove, for a fixed $z$ of norm at most
$R$, that
\begin{equation}
    \Prob\{|\ip{u}{z}|\geq\tau\}
    \leq2\exp\!\left(-\frac{(k-1)\tau^2}{2R^2}\right).
    \label{eq:single-cap}
\end{equation}
Write $p=P_Wz$ and $r=\norm{p}$. If $r<\tau$, the probability is zero.
If $r=\tau$, equality could occur only at the two sphere points parallel
to $p$, a set of probability zero because $k\geq2$.
It therefore suffices to consider $r>\tau$.

Choose orthonormal coordinates in $W$ with $p=re_1$ and represent
$u=g/\norm{g}$, where $g_1,\ldots,g_k$ are independent standard Gaussian
variables. Put
\[
    \alpha:=\tau/r\in(0,1),\qquad S:=\sum_{i=2}^k g_i^2.
\]
The cap event is equivalent to
\[
    |g_1|\geq\frac{\alpha}{\sqrt{1-\alpha^2}}\sqrt S.
\]
For a standard Gaussian $G$, the exponential Markov inequality gives
$\Prob(G\geq t)\leq e^{-t^2/2}$ for $t\geq0$; applying it to $G$ and
$-G$ yields $\Prob(|G|\geq t)\leq2e^{-t^2/2}$.
Conditioning on $S$, using independence, and then integrating gives
\begin{align*}
    \Prob\{|\ip{u}{z}|\geq\tau\}
    &\leq2\E\exp\!\left(-\frac{\alpha^2}{2(1-\alpha^2)}S\right)\\
    &=2(1-\alpha^2)^{(k-1)/2}\\
    &\leq2\exp\!\left(-\frac{k-1}{2}\alpha^2\right).
\end{align*}
The equality follows from
$\E e^{-tG^2/2}=(1+t)^{-1/2}$ for $t\geq0$, obtained by integrating the
Gaussian density, and taking a product over $k-1$ variables.
The last inequality uses $\log(1-s)\leq-s$ for $0\leq s<1$.
Since $r\leq R$, we have $\alpha\geq\tau/R$, proving \eqref{eq:single-cap}.

A union bound over the $M$ fixed vectors bounds the probability of any
failure in \eqref{eq:cap-conclusion} by the left side of
\eqref{eq:cap-condition}. That number is strictly below one, so a vector
satisfying all the strict inequalities exists.
\end{proof}

\subsection{Delayed directions and the final output}

Fix $d\geq8$ and an integer $T\geq1$, and define
\begin{equation}
    b:=\left\lfloor\frac d4\right\rfloor,\qquad
    q:=\left\lceil\frac Tb\right\rceil,\qquad n:=q+1.
    \label{eq:blocks}
\end{equation}
The radius $R=D\sqrt n$ and the map $\rho$ are those in
\eqref{eq:radial-constants}--\eqref{eq:rho}. Assume for the following lemma
that
\begin{equation}
    n\leq\frac d4,\qquad
    \frac{d}{32D^2n}\geq4\log(ed),\qquad T\leq d^2.
    \label{eq:frame-hypotheses}
\end{equation}

\begin{lemma}[Delayed-frame completion]
\label{lem:completion}
Fix positive numbers $\mu,s$ before the interaction, and let
$A\in\DetVal_T$. Under \eqref{eq:frame-hypotheses}, there are replies
$r_1,\ldots,r_T$ and an orthonormal frame $U=(u_1,\ldots,u_n)$ such that
\begin{equation}
    g_U(x):=\mu F_U(x/s)
    \label{eq:gU}
\end{equation}
has the following properties. If $x_t$ and $x_{\mathrm{out}}$ are the
queries and output generated by the reply sequence $(r_1,\ldots,r_T)$,
then
\begin{equation}
    g_U(x_t)=r_t\quad(1\leq t\leq T),\qquad
    \ip{u_n}{\rho(x_{\mathrm{out}}/s)}=0.
    \label{eq:completion-properties}
\end{equation}
In particular, the single fixed objective $g_U$ produces exactly these
queries, replies, and output when used as the oracle for $A$.
\end{lemma}

\begin{proof}
Partition the query indices into the consecutive, nonempty blocks
\[
    B_j:=\{(j-1)b+1,\ldots,\min(jb,T)\},\qquad 1\leq j\leq q.
\]
We first construct $u_1,\ldots,u_q$ one block at a time.

\paragraph{Generating a block before choosing its direction.}
Suppose $u_1,\ldots,u_{j-1}$ have been selected. For each query $x_t$ in
block $B_j$, set
\begin{equation}
    w_t:=\rho(x_t/s)
    \label{eq:wt}
\end{equation}
and return
\begin{equation}
    r_t:=\mu\left[
      h_n\bigl(\ip{u_1}{w_t},\ldots,\ip{u_{j-1}}{w_t},0,\ldots,0\bigr)
         +\eta\norm{x_t/s}^2\right].
    \label{eq:online-reply}
\end{equation}
There are exactly $n$ arguments in this evaluation of $h_n$; for $j=1$
they are all zero. The reply uses no unselected direction.
By recursion through the deterministic maps of the policy, every query
and reply within this block is fixed before $u_j$ is selected.

\paragraph{Choosing the next direction.}
After the block has been generated, let
\begin{equation}
    W_j:=\Span\bigl(\{u_1,\ldots,u_{j-1}\}
                     \cup\{w_t:t\in B_j\}\bigr)^\perp.
    \label{eq:Wj}
\end{equation}
Because $j-1\leq q-1=n-2$, $|B_j|\leq b\leq d/4$, and $n\leq d/4$,
\begin{equation}
    k:=\dim W_j\geq d-(n-2)-b\geq d/2.
    \label{eq:dimension}
\end{equation}
In particular, $k\geq2$ and $k-1\geq d/4$.
Every historical point has $\norm{w_t}<R$.
Apply \cref{lem:caps} in $W_j$ with $\tau=1/2$ to the points from strictly
earlier blocks. There are at most $T\leq d^2$ such points, and
\begin{equation}
    \frac{(k-1)\tau^2}{2R^2}
    =\frac{k-1}{8D^2n}\geq\frac{d}{32D^2n}
    \geq4\log(ed).
    \label{eq:cap-exponent}
\end{equation}
Consequently the union bound is at most
\[
    2d^2e^{-4\log(ed)}=2e^{-4}d^{-2}<1.
\]
For the first block the list of earlier points is empty, which is also
covered by the lemma. We may therefore choose a unit vector $u_j\in W_j$
satisfying
\begin{equation}
\begin{aligned}
    \ip{u_j}{w_t}&=0 &&(t\in B_j),\\
    |\ip{u_j}{w_t}|&<1/2 &&(t<\min B_j).
\end{aligned}
\label{eq:direction-invariant}
\end{equation}
Membership in $W_j$ ensures that it is orthogonal to all previously chosen
directions. This completes the induction over $j=1,\ldots,q$.

\paragraph{Choosing the final direction after the output.}
All $T$ replies now have been fixed. Let
\[
    x_{\mathrm{out}}:=\pi_{\mathrm{out}}(r_1,\ldots,r_T),\qquad
    w_{\mathrm{out}}:=\rho(x_{\mathrm{out}}/s),
\]
and define
\begin{equation}
    W_{\mathrm{out}}:=\Span\{u_1,\ldots,u_q,w_{\mathrm{out}}\}^\perp.
    \label{eq:Wout}
\end{equation}
Its dimension is at least $d-q-1=d-n\geq3d/4$.
The same cap estimate, now applied to all $w_1,\ldots,w_T$, provides a
unit vector $u_n\in W_{\mathrm{out}}$ with
\begin{equation}
    \ip{u_n}{w_{\mathrm{out}}}=0,\qquad
    |\ip{u_n}{w_t}|<1/2\quad(1\leq t\leq T).
    \label{eq:sentinel}
\end{equation}
Together with $u_1,\ldots,u_q$, this gives an orthonormal frame $U$.

\paragraph{Preservation of every exact value.}
Fix the completed frame and a historical query $t\in B_j$.
By \eqref{eq:direction-invariant}, its $j$th coordinate is exactly zero:
$\ip{u_j}{w_t}=0$. For each $j<k\leq q$, query $t$ belongs to an earlier
block than $B_k$, so the second line of \eqref{eq:direction-invariant}
gives $|\ip{u_k}{w_t}|<1/2$. For $k=n$, the same inequality follows from
\eqref{eq:sentinel}. Applying \eqref{eq:shield} therefore yields
\[
    h_n(U^\top w_t)
    =h_n\bigl(\ip{u_1}{w_t},\ldots,\ip{u_{j-1}}{w_t},0,\ldots,0\bigr).
\]
The quadratic term and the scaling were fixed throughout the interaction.
Equations \eqref{eq:FU}, \eqref{eq:gU}, and \eqref{eq:online-reply} now show
that $g_U(x_t)=r_t$ for every $t$.
This identity also covers repeated query points, which necessarily receive
the same value from the completed function.

\paragraph{Replay against a single fixed objective.}
The first query of $A$ is independent of any reply and hence agrees with
the simulated first query; the completed objective returns $r_1$ there.
If the first $t-1$ replies agree, determinism forces the same $t$th query,
and the identity $g_U(x_t)=r_t$ gives the same next reply.
Induction proves equality of the entire transcript, and the deterministic
output map then returns the same $x_{\mathrm{out}}$.
The second identity in \eqref{eq:completion-properties} is precisely
\eqref{eq:sentinel}.
\end{proof}

\begin{remark}[No randomization or continuity assumption on the policy]
The Gaussian variables in \cref{lem:caps} are a device for proving the
existence of a geometric direction. At each use of that lemma the relevant
subspace and all tested points are fixed. No probability distribution is
placed on the policy, and no measurability or continuity of its maps is
needed. Conversely, a frame selected against a deterministic policy in
this manner is not a policy-independent distribution of hard functions.
\end{remark}

\section{Scaling and proof of the dimension-dependent bound}
\label{sec:lower}

\subsection{Reduction to normalized parameters}

\begin{lemma}[Exact-value scaling]
\label{lem:scaling}
A deterministic lower bound with normalized parameters
\[
    L=1,\qquad \epsilon=1,\qquad\Delta=a,\qquad x_0=0
\]
transfers, with the same query budget, to arbitrary parameters satisfying
$a=L\Delta/\epsilon^2$. Exact transcript consistency, the exact initial
gap, and strict failure at the output are preserved.
\end{lemma}

\begin{proof}
Fix a policy $A=(\pi_t,\pi_{\mathrm{out}})$ for the original parameters.
Define a normalized policy $B$ by
\begin{equation}
    z_t=\frac L\epsilon\left[
       \pi_t\left(\frac{\epsilon^2}{L}w_1,\ldots,
                  \frac{\epsilon^2}{L}w_{t-1}\right)-x_0\right]
    \label{eq:scaled-policy}
\end{equation}
on a normalized history $(w_1,\ldots,w_{t-1})$, and use the analogous
formula for its output $Z_{\mathrm{out}}$.
Suppose a fixed normalized objective $g$ reproduces this policy's history,
has $\Lip(\nabla g)\leq1$ and $g(0)-\inf g=a$, and satisfies
$\norm{\nabla g(Z_{\mathrm{out}})}>1$.
Define
\begin{equation}
    f(x):=\frac{\epsilon^2}{L}
             g\!\left(\frac L\epsilon(x-x_0)\right).
    \label{eq:scaled-function}
\end{equation}
Then
\begin{equation}
    \nabla f(x)=\epsilon\nabla g\!\left(\frac L\epsilon(x-x_0)\right),
    \qquad
    \nabla^2 f(x)=L\nabla^2g\!\left(\frac L\epsilon(x-x_0)\right).
    \label{eq:scaled-derivatives}
\end{equation}
It follows that $f\in C^\infty$, $\Lip(\nabla f)\leq L$, and
\[
    f(x_0)-\inf f=\frac{\epsilon^2}{L}(g(0)-\inf g)
                =\frac{\epsilon^2}{L}a=\Delta.
\]
Lower boundedness and, when present, coercivity are preserved.

To check the actual transcript, suppose the first $t-1$ replies of $A$
are $(\epsilon^2/L)w_1,\ldots,(\epsilon^2/L)w_{t-1}$.
Equation \eqref{eq:scaled-policy} then gives
$x_t=x_0+(\epsilon/L)z_t$, and \eqref{eq:scaled-function} gives
$f(x_t)=(\epsilon^2/L)g(z_t)=(\epsilon^2/L)w_t$.
This proves correspondence of the replies by induction, starting with the
empty history. The output corresponds by the same definition, so
\[
    \norm{\nabla f(X_{\mathrm{out}}^A(f))}
    =\epsilon\norm{\nabla g(Z_{\mathrm{out}})}>\epsilon.
\]
\end{proof}

\subsection{An explicit normalized lower bound}

The constants $C_h,G,H,D,\Lambda$ are defined in
\eqref{eq:chain-constants}, \eqref{eq:radial-constants}, and
\eqref{eq:Lambda}; $\gamma=3/4$ is the constant in
\cref{lem:certificate}. Fix
\begin{equation}
    c:=\min\left\{\frac1{48},\frac1{1536D^2},
                      \frac{\gamma^2}{24C_h\Lambda}\right\}>0.
    \label{eq:c}
\end{equation}
These constants are universal and are not optimized.

\begin{theorem}[Normalized lower bound]
\label{thm:normalized}
Let $d\geq8$, $a>1/2$, and
\begin{equation}
    m_\star:=\min\left\{a,\frac{d}{\log(ed)}\right\}.
    \label{eq:mstar}
\end{equation}
For every integer $T<cdm_\star$ and every $A\in\DetVal_T$, there is one
fixed coercive $g\in C^\infty(\R^d)$ such that
\begin{equation}
    \Lip(\nabla g)\leq1,\qquad g(0)-\inf g=a,\qquad
    \norm{\nabla g(X_{\mathrm{out}}^A(g))}>1.
    \label{eq:normalized-result}
\end{equation}
The function exactly realizes the complete charged transcript of the
policy.
\end{theorem}

\begin{proof}
Fix such a budget and policy. If $T\leq d$, apply \cref{prop:quadratic} with
$L=\epsilon=1$, $\Delta=a$, and $x_0=0$. Its proof supplies a coercive
quadratic satisfying all the assertions. We may therefore assume $T>d$.

\paragraph{Chain length and dimension requirements.}
Define $b,q,n$ by \eqref{eq:blocks}. Since $d\geq8$,
$b=\lfloor d/4\rfloor\geq d/8$. Since $T/d\geq1$,
\begin{equation}
    n=\left\lceil\frac Tb\right\rceil+1
      \leq\frac{8T}{d}+2
      \leq\frac{10T}{d}
      <\frac{12T}{d}
      <12c m_\star.
    \label{eq:length}
\end{equation}
In particular,
\begin{equation}
    n<12ca,\qquad n<\frac{12cd}{\log(ed)}.
    \label{eq:length-both}
\end{equation}
The choice $c\leq1/48$ and the inequality $\log(ed)\geq1$ imply
$n\leq d/4$. Also
$T<cdm_\star\leq cd^2/\log(ed)<d^2$.
For the cap condition,
\begin{equation}
    \frac{d}{32D^2n}
       >\frac{\log(ed)}{384D^2c}\geq4\log(ed),
    \label{eq:cap-accounting}
\end{equation}
where the last step uses $c\leq1/(1536D^2)$.
Thus every hypothesis in \eqref{eq:frame-hypotheses} holds.

\paragraph{Fixing the exact-gap scaling before interaction.}
For this $n$ and $d$, let $\delta_{n,d}$ be the orientation-independent gap
from \cref{lem:family}. Define
\begin{equation}
    \mu:=\frac{a}{\delta_{n,d}},\qquad s:=\sqrt{\mu\Lambda}.
    \label{eq:mu-s}
\end{equation}
Both numbers are finite and strictly positive by \eqref{eq:gap-bounds}.
They depend only on $a,n,d$ and the fixed analytic family, not on the
orientation eventually chosen against the policy. They may therefore be
fixed before any replies are generated.
Furthermore, using \eqref{eq:gap-bounds}, \eqref{eq:length-both}, and
\eqref{eq:c},
\begin{equation}
    \mu\geq\frac{a}{C_hn}
          >\frac{1}{12C_hc}
          \geq\frac{2\Lambda}{\gamma^2}.
    \label{eq:amplitude}
\end{equation}

\paragraph{Completing the objective and checking its properties.}
Apply \cref{lem:completion} with these fixed $\mu,s$.
It produces an orthonormal frame $U$ and the fixed objective
$g_U(x)=\mu F_U(x/s)$ that reproduces the entire transcript.
The function is $C^\infty$ and coercive. Its smoothness satisfies
\begin{equation}
    \Lip(\nabla g_U)\leq\frac{\mu}{s^2}\Lambda=1,
    \label{eq:normalized-smoothness}
\end{equation}
and the bijective change of variables $x\mapsto x/s$ gives the exact gap
\begin{equation}
    g_U(0)-\inf g_U=\mu\delta_{n,d}=a.
    \label{eq:normalized-gap}
\end{equation}
At the policy's actual output, \cref{lem:completion} gives
$\ip{u_n}{\rho(x_{\mathrm{out}}/s)}=0$.
By \cref{lem:certificate} and \eqref{eq:amplitude},
\begin{equation}
    \norm{\nabla g_U(x_{\mathrm{out}})}
       =\frac{\mu}{s}\norm{\nabla F_U(x_{\mathrm{out}}/s)}
       >\gamma\sqrt{\frac{\mu}{\Lambda}}>
         \sqrt2>1.
    \label{eq:normalized-failure}
\end{equation}
This proves every assertion of \eqref{eq:normalized-result} in the case
$T>d$, and the two cases exhaust all admissible integer budgets.
\end{proof}

\subsection{Completion of the main theorem}
\label{sec:main-proof}

\begin{proof}[Proof of \cref{thm:main}]
The implication $a\leq1/2\Rightarrow\Nval(d,L,\Delta)=0$ follows
from \eqref{eq:start}. If $a>1/2$, \cref{prop:quadratic} gives
$\Nval(d,L,\Delta)\geq d+1\geq1$, proving the converse and the lower part
of \eqref{eq:basic-bounds}. The upper part follows from
\cref{prop:upper}. Equation \eqref{eq:transition} proves the claimed
uniform $\Theta(d)$ statement.

For $d\geq8$ and $a>1/2$, fix any integer
$T<cd\min\{a,d/\log(ed)\}$ and any original policy.
Normalize it using \cref{lem:scaling}, defeat the normalized policy by
\cref{thm:normalized}, and scale the resulting objective back.
This gives one fixed admissible objective for which the actual output
fails. Thus no integer query budget below the right-hand side of
\eqref{eq:main-lower} can suffice. Since a successful finite budget exists
by the upper bound, the definition \eqref{eq:complexity} implies
\eqref{eq:main-lower} without any loss from integer rounding.

In the range $1/2<a\leq d/\log(ed)$, the minimum equals $a$, while
$2d\leq4da$. Hence \eqref{eq:basic-bounds} gives
$\Nval(d,L,\Delta)\leq8da+2d\leq12da$, proving \eqref{eq:matching}.
Both lower-bound constructions are coercive and explicitly reproduce the
complete exact transcript, establishing the final statement of the theorem.
\end{proof}


\section{Discussion and limitations}
\label{sec:scope}

The main theorem gives a complete deterministic rate, up to universal
constants, in the range $d\geq8$ and $1/2<a\leq d/\log(ed)$.
Within this range the factor $d$ in the coordinate finite-difference
method is unavoidable even when algorithms can query arbitrarily distant
points and return a point outside their query set. Exact initial-gap
normalization and a final direction selected against the output are part
of the proof, rather than additional restrictions on the algorithm.

\paragraph{The high-accuracy range.}
For $d\geq8$ and $a>d/\log(ed)$, the present results give only
\begin{equation}
 \max\!\left\{d+1,\frac{cd^2}{\log(ed)}\right\}
 \leq\Nval(d,L,\Delta)
 \leq(d+1)\lceil4a\rceil.
 \label{eq:remaining-range}
\end{equation}
The logarithm enters when one direction must have small projections onto
all earlier query points. With compression radius of order $\sqrt n$,
the cap estimate requires $n\log(ed)$ to be at most a constant times $d$.
This is a limitation of the construction. The theorem supplies neither a
matching upper bound for the saturated lower-bound branch nor a proof
that the logarithm is necessary. In dimensions $1\leq d<8$, the elementary
bounds still give $\Theta(d)$ for $1/2<a\leq1$; they do not determine the
dependence on $a$ at higher accuracy.

\paragraph{Randomized policies.}
The directions in our proof are chosen against a fixed deterministic
policy, with the last direction chosen after its output is known. This
does not produce a single instance or distribution defeating a randomized
algorithm. The randomized lower bound of \citet{XY2026} applies under its
own hypotheses, as stated in \Cref{tab:comparison}; extending the present
exact-gap construction to an analogous randomized statement is a separate
question.

\paragraph{Oracle conventions and regularity.}
The $d+1$ elementary lower bound uses the fact that the unknown objective
may have any additive constant and that $f(x_0)$ is not supplied for free.
The exact quadratic calculation in \Cref{cor:quadratic-exact} makes this
convention visible. The main multiplicative rate does not assert
arithmetic complexity, robustness to rounding, or hardness under a
stochastic observation model. The hard objectives are $C^\infty$ and
coercive, but the class controls only their first-order smoothness.

The proof isolates the two ingredients responsible for the multiplicative
cost: exact invariance of historical scalar values and a gradient
obstruction at the terminal output. The former allows a full block of
queries to precede the selection of a direction; the latter ensures that
an arbitrary output cannot bypass the hidden chain. Together with
orientation-independent gap normalization, they yield a self-contained
deterministic lower bound with the same dimension and accuracy order as
finite-difference descent in the matched range.

\begin{thebibliography}{99}

\bibitem[Carmon et~al.(2020)Carmon, Duchi, Hinder, and Sidford]{CDHS2020}
Y.~Carmon, J.~C. Duchi, O.~Hinder, and A.~Sidford.
\newblock Lower bounds for finding stationary points I.
\newblock \emph{Mathematical Programming}, 184:71--120, 2020.
\newblock \href{https://doi.org/10.1007/s10107-019-01406-y}{doi:10.1007/s10107-019-01406-y}.

\bibitem[Wu et~al.(2026a)Wu, Zhang, Zhang, Li, Zheng, Fang, Li, and Lin]{WuEtAl2026Convex}
W.~Wu, H.~Zhang, C.~Zhang, Y.~Li, C.~Zheng, C.~Fang, H.~Li, and Z.~Lin.
\newblock Near-optimal deterministic exact-value complexity for smooth convex optimization.
\newblock \emph{arXiv preprint arXiv:2609.18230}, 2026.
\newblock \url{https://arxiv.org/abs/2609.18230}.

\bibitem[Wu et~al.(2026b)Wu, Zhang, Zhang, Li, Zheng, Fang, Li, and Lin]{WuEtAl2026SC}
W.~Wu, H.~Zhang, C.~Zhang, Y.~Li, C.~Zheng, C.~Fang, H.~Li, and Z.~Lin.
\newblock Near-optimal exact-value zeroth-order complexity for smooth strongly convex optimization.
\newblock \emph{arXiv preprint arXiv:2609.18728}, 2026.
\newblock \url{https://arxiv.org/abs/2609.18728}.

\bibitem[Peng et~al.(2026)Peng, Tang, and Jia]{PTJ2026}
Y.~Peng, Z.~Tang, and W.~Jia.
\newblock Lower bounds for smooth nonconvex zeroth-order optimization:
The price of determinism under low curvature trace.
\newblock \emph{arXiv preprint arXiv:2610.27550}, 2026.
\newblock \url{https://arxiv.org/abs/2610.27550}.

\bibitem[Nesterov and Spokoiny(2017)]{NS2017}
Y.~Nesterov and V.~Spokoiny.
\newblock Random gradient-free minimization of convex functions.
\newblock \emph{Foundations of Computational Mathematics}, 17:527--566, 2017.
\newblock \href{https://doi.org/10.1007/s10208-015-9296-2}{doi:10.1007/s10208-015-9296-2}.

\bibitem[Ghadimi and Lan(2013)]{GL2013}
S.~Ghadimi and G.~Lan.
\newblock Stochastic first- and zeroth-order methods for nonconvex stochastic
programming.
\newblock \emph{SIAM Journal on Optimization}, 23(4):2341--2368, 2013.
\newblock \href{https://doi.org/10.1137/120880811}{doi:10.1137/120880811}.

\bibitem[Xia and Yuan(2026)]{XY2026}
Y.~Xia and X.-T. Yuan.
\newblock On global joint lower bounds in zeroth-order nonconvex optimization.
\newblock \emph{arXiv preprint arXiv:2610.19652v2}, 2026.
\newblock \url{https://arxiv.org/abs/2610.19652v2}.

\bibitem[Conn et~al.(2009)Conn, Scheinberg, and Vicente]{CSV2009}
A.~R. Conn, K.~Scheinberg, and L.~N. Vicente.
\newblock \emph{Introduction to Derivative-Free Optimization}.
\newblock MPS-SIAM Series on Optimization. SIAM, Philadelphia, 2009.
\newblock \url{https://www.mat.uc.pt/~lnv/idfo/}.

\bibitem[Larson et~al.(2019)Larson, Menickelly, and Wild]{LMW2019}
J.~Larson, M.~Menickelly, and S.~M. Wild.
\newblock Derivative-free optimization methods.
\newblock \emph{Acta Numerica}, 28:287--404, 2019.
\newblock \href{https://doi.org/10.1017/S0962492919000060}{doi:10.1017/S0962492919000060}.

\bibitem[Conn et~al.(2009)Conn, Scheinberg, and Vicente]{CSV2009TR}
A.~R. Conn, K.~Scheinberg, and L.~N. Vicente.
\newblock Global convergence of general derivative-free trust-region algorithms to first- and second-order critical points.
\newblock \emph{SIAM Journal on Optimization}, 20(1):387--415, 2009.
\newblock \href{https://doi.org/10.1137/060673424}{doi:10.1137/060673424}.

\bibitem[Garmanjani et~al.(2016)Garmanjani, J\'udice, and Vicente]{GJV2016}
R.~Garmanjani, D.~J\'udice, and L.~N. Vicente.
\newblock Trust-region methods without using derivatives: Worst case complexity and the nonsmooth case.
\newblock \emph{SIAM Journal on Optimization}, 26(4):1987--2011, 2016.
\newblock \href{https://doi.org/10.1137/151005683}{doi:10.1137/151005683}.

\bibitem[Vicente(2013)]{Vicente2013}
L.~N. Vicente.
\newblock Worst case complexity of direct search.
\newblock \emph{EURO Journal on Computational Optimization}, 1:143--153, 2013.
\newblock \href{https://doi.org/10.1007/s13675-012-0003-7}{doi:10.1007/s13675-012-0003-7}.

\bibitem[Cartis et~al.(2012)Cartis, Gould, and Toint]{CGT2012}
C.~Cartis, N.~I.~M. Gould, and P.~L. Toint.
\newblock On the oracle complexity of first-order and derivative-free algorithms for smooth nonconvex minimization.
\newblock \emph{SIAM Journal on Optimization}, 22(1):66--86, 2012.
\newblock \href{https://doi.org/10.1137/100812276}{doi:10.1137/100812276}.

\bibitem[Dzahini et~al.(2025)Dzahini, Rinaldi, Royer, and Zeffiro]{DRRZ2025}
K.~J. Dzahini, F.~Rinaldi, C.~W. Royer, and D.~Zeffiro.
\newblock Direct-search methods in the year 2025: Theoretical guarantees and algorithmic paradigms.
\newblock \emph{EURO Journal on Computational Optimization}, 13:100110, 2025.
\newblock \href{https://doi.org/10.1016/j.ejco.2025.100110}{doi:10.1016/j.ejco.2025.100110}.

\bibitem[Gratton et~al.(2015)Gratton, Royer, Vicente, and Zhang]{GRVZ2015}
S.~Gratton, C.~W. Royer, L.~N. Vicente, and Z.~Zhang.
\newblock Direct search based on probabilistic descent.
\newblock \emph{SIAM Journal on Optimization}, 25(3):1515--1541, 2015.
\newblock \href{https://doi.org/10.1137/140961602}{doi:10.1137/140961602}.

\bibitem[Vavasis(1993)]{Vavasis1993}
S.~A. Vavasis.
\newblock Black-box complexity of local minimization.
\newblock \emph{SIAM Journal on Optimization}, 3(1):60--80, 1993.
\newblock \href{https://doi.org/10.1137/0803004}{doi:10.1137/0803004}.

\bibitem[Kerger(2026)]{Kerger2026}
P.~Kerger.
\newblock Closing the oracle-complexity gap in derivative-free convex optimization: A near-quadratic lower bound from exact function values.
\newblock \emph{arXiv preprint arXiv:2607.13335}, 2026.
\newblock \url{https://arxiv.org/abs/2607.13335}.

\bibitem[Shamir(2013)]{Shamir2013}
O.~Shamir.
\newblock On the complexity of bandit and derivative-free stochastic convex optimization.
\newblock In \emph{Proceedings of the 26th Annual Conference on Learning Theory}, volume 30, pages 3--24. PMLR, 2013.
\newblock \url{https://proceedings.mlr.press/v30/Shamir13.html}.

\bibitem[Nesterov and Polyak(2006)]{NP2006}
Y.~Nesterov and B.~T. Polyak.
\newblock Cubic regularization of Newton method and its global performance.
\newblock \emph{Mathematical Programming}, 108:177--205, 2006.
\newblock \href{https://doi.org/10.1007/s10107-006-0706-8}{doi:10.1007/s10107-006-0706-8}.

\bibitem[Carmon et~al.(2017)Carmon, Duchi, Hinder, and Sidford]{CDHS2017}
Y.~Carmon, J.~C. Duchi, O.~Hinder, and A.~Sidford.
\newblock ``Convex until proven guilty'': Dimension-free acceleration of gradient descent on non-convex functions.
\newblock In \emph{Proceedings of the 34th International Conference on Machine Learning}, volume 70, pages 654--663. PMLR, 2017.
\newblock \url{https://proceedings.mlr.press/v70/carmon17a.html}.

\bibitem[Carmon et~al.(2021)Carmon, Duchi, Hinder, and Sidford]{CDHS2021}
Y.~Carmon, J.~C. Duchi, O.~Hinder, and A.~Sidford.
\newblock Lower bounds for finding stationary points II: First-order methods.
\newblock \emph{Mathematical Programming}, 185:315--355, 2021.
\newblock \href{https://doi.org/10.1007/s10107-019-01431-x}{doi:10.1007/s10107-019-01431-x}.

\bibitem[Doikov and Grapiglia(2025)]{DG2025}
N.~Doikov and G.~N. Grapiglia.
\newblock First and zeroth-order implementations of the regularized Newton method with lazy approximated Hessians.
\newblock \emph{Journal of Scientific Computing}, 103:32, 2025.
\newblock \href{https://doi.org/10.1007/s10915-025-02842-z}{doi:10.1007/s10915-025-02842-z}.

\bibitem[Hinder(2018)]{Hinder2018}
O.~Hinder.
\newblock Cutting plane methods can be extended into nonconvex optimization.
\newblock In \emph{Proceedings of the 31st Conference on Learning Theory}, volume 75, pages 1451--1454. PMLR, 2018.
\newblock \url{https://proceedings.mlr.press/v75/hinder18a.html}.

\bibitem[Chewi et~al.(2023)Chewi, Bubeck, and Salim]{CBS2023}
S.~Chewi, S.~Bubeck, and A.~Salim.
\newblock On the complexity of finding stationary points of smooth functions in one dimension.
\newblock In \emph{Proceedings of the 34th International Conference on Algorithmic Learning Theory}, volume 201, pages 358--374. PMLR, 2023.
\newblock \url{https://proceedings.mlr.press/v201/chewi23a.html}.

\end{thebibliography}

\end{document}
